% GENERATED FILE. Edit canonical lessons, then run npm run generate:books.
% canonical-source-hash: fnv1a64:e44dd8900c81b23e
% canonical-lessons: GU-C507-parovvu, GU-C507-moti, GU-C507-sandhvu, GU-C507-thigdu, GU-C507-bharatkam

\chapter{Needle and Thread}
\label{ch:gu-needle-and-thread}
\hlchaptermodality{507}

\begin{chapteropening}
\textbf{By the end of this chapter:} \emph{I can talk about threading a needle, mending and embroidery.}

Complete the last lesson of chapter 507: I can talk about threading a needle, mending and embroidery.
\end{chapteropening}

\section[\texorpdfstring{\gu{પરોવવું} (parovvũ)}{પરોવવું (parovvũ)}]{\gu{પરોવવું} (parovvũ) — to thread (a needle)}

\label{lesson:GU-C507-parovvu}



\begin{quote}
Before the new one: say the Gujarati for to stretch tight, then the Gujarati for to wind round.
\end{quote}

\subsection*{You'll want to know: \gu{પરોવવું}}
\textbf{\gu{પરોવવું}} — \emph{parovvũ} — \textquotedblleft{}to thread (a needle)\textquotedblright{}.

\begin{quote}
\gu{દાદીએ} \gu{ચશ્માં} \gu{પહેરીને} \gu{સોયમાં} \gu{દોરો} \gu{પરોવ્યો}.
\end{quote}

\begin{quote}
\emph{dādīe caśmā̃ paherīne soymā̃ doro parovyo.}
\end{quote}

\begin{quote}
Grandmother put on her glasses and threaded the needle.
\end{quote}

Past, with \gu{તેણે} (he or she …): \gu{તેણે} \gu{દોરો} \gu{પરોવ્યો}, \emph{teṇe doro parovyo}.

\begin{cousinweb}
\textbf{\gu{પરોવવું}} (\emph{parovvũ}) is inherited from Sanskrit \emph{právayati}, \textquotedblleft{}weaves on\textquotedblright{}, and CDIAL lists the Gujarati verb, \textquotedblleft{}to thread a needle, string pearls\textquotedblright{} (CDIAL 8785). What goes through the needle is the object: \textbf{\gu{દોરો} \gu{પરોવ્યો}}.
\end{cousinweb}

\begin{grammarlens}[title={What you've built}]
One more word for describing what you lived and felt.
\end{grammarlens}

\subsection*{Your turn}
\begin{itemize}
\raggedright
  \item \emph{Say it:} \emph{parovvũ}
  \item \emph{Say it:} \emph{parovvũ}, once more
  \item \emph{Say it:} \emph{dādīe caśmā̃ paherīne soymā̃ doro parovyo.} — Grandmother put on her glasses and threaded the needle.
  \item \emph{Recall:} say the Gujarati for to stretch tight, then the Gujarati for to wind round, then say \emph{parovvũ} again
\end{itemize}

\begin{tcolorbox}[breakable,colback=teal!4,colframe=teal!35!black,title={Before you move on}]
What does \gu{પરોવવું} mean? (\textquotedblleft{}To thread (a needle)\textquotedblright{}.) Say it once more.
\end{tcolorbox}
\section[\texorpdfstring{\gu{મોતી} (motī)}{મોતી (motī)}]{\gu{મોતી} (motī) — a pearl}

\label{lesson:GU-C507-moti}



\begin{quote}
Before the new one: say the Gujarati for to wind round, then the Gujarati for to thread (a needle).
\end{quote}

\subsection*{You'll want to know: \gu{મોતી}}
\textbf{\gu{મોતી}} — \emph{motī} — \textquotedblleft{}a pearl\textquotedblright{}.

\begin{quote}
\gu{તેણે} \gu{દોરામાં} \gu{મોતી} \gu{પરોવ્યાં}.
\end{quote}

\begin{quote}
\emph{teṇe dorāmā̃ motī parovyā̃.}
\end{quote}

\begin{quote}
She strung pearls on a thread.
\end{quote}

A neuter noun.

\begin{cousinweb}
\textbf{\gu{મોતી}} (\emph{motī}) is a tadbhava, inherited from Sanskrit \emph{mauktika}, \textquotedblleft{}a pearl\textquotedblright{}, and CDIAL lists the Gujarati word (CDIAL 10365). It is neuter, and the plural looks the same: \textbf{\gu{મોતી} \gu{પરોવ્યાં}}.
\end{cousinweb}

\begin{grammarlens}[title={What you've built}]
One more word for describing what you lived and felt.
\end{grammarlens}

\subsection*{Your turn}
\begin{itemize}
\raggedright
  \item \emph{Say it:} \emph{motī}
  \item \emph{Say it:} \emph{motī}, once more
  \item \emph{Say it:} \emph{teṇe dorāmā̃ motī parovyā̃.} — She strung pearls on a thread.
  \item \emph{Recall:} say the Gujarati for to wind round, then the Gujarati for to thread (a needle), then say \emph{motī} again
\end{itemize}

\begin{tcolorbox}[breakable,colback=teal!4,colframe=teal!35!black,title={Before you move on}]
What does \gu{મોતી} mean? (\textquotedblleft{}A pearl\textquotedblright{}.) Say it once more.
\end{tcolorbox}
\section[\texorpdfstring{\gu{સાંધવું} (sā̃dhvũ)}{સાંધવું (sā̃dhvũ)}]{\gu{સાંધવું} (sā̃dhvũ) — to mend (a tear)}

\label{lesson:GU-C507-sandhvu}



\begin{quote}
Before the new one: say the Gujarati for to thread (a needle), then the Gujarati for a pearl.
\end{quote}

\subsection*{You'll want to know: \gu{સાંધવું}}
\textbf{\gu{સાંધવું}} — \emph{sā̃dhvũ} — \textquotedblleft{}to mend (a tear)\textquotedblright{}.

\begin{quote}
\gu{દરજીએ} \gu{મારું} \gu{જૂનું} \gu{ખમીસ} \gu{સાંધ્યું}.
\end{quote}

\begin{quote}
\emph{darjīe mārũ jūnũ khamīs sā̃dhyũ.}
\end{quote}

\begin{quote}
The tailor mended my old shirt.
\end{quote}

Past, with \gu{તેણે} (he or she …): \gu{તેણે} \gu{ખમીસ} \gu{સાંધ્યું}, \emph{teṇe khamīs sā̃dhyũ}.

\begin{cousinweb}
\textbf{\gu{સાંધવું}} (\emph{sā̃dhvũ}) is inherited from Sanskrit \emph{sáṁdadhāti}, \textquotedblleft{}places on, puts together\textquotedblright{}, and CDIAL lists the Gujarati verb, \textquotedblleft{}to unite, attach\textquotedblright{} (CDIAL 12898).
\end{cousinweb}

\begin{grammarlens}[title={What you've built}]
One more word for describing what you lived and felt.
\end{grammarlens}

\subsection*{Your turn}
\begin{itemize}
\raggedright
  \item \emph{Say it:} \emph{sā̃dhvũ}
  \item \emph{Say it:} \emph{sā̃dhvũ}, once more
  \item \emph{Say it:} \emph{darjīe mārũ jūnũ khamīs sā̃dhyũ.} — The tailor mended my old shirt.
  \item \emph{Recall:} say the Gujarati for to thread (a needle), then the Gujarati for a pearl, then say \emph{sā̃dhvũ} again
\end{itemize}

\begin{tcolorbox}[breakable,colback=teal!4,colframe=teal!35!black,title={Before you move on}]
What does \gu{સાંધવું} mean? (\textquotedblleft{}To mend (a tear)\textquotedblright{}.) Say it once more.
\end{tcolorbox}
\section[\texorpdfstring{\gu{થીંગડું} (thī̃gḍũ)}{થીંગડું (thī̃gḍũ)}]{\gu{થીંગડું} (thī̃gḍũ) — a patch (on cloth)}

\label{lesson:GU-C507-thigdu}



\begin{quote}
Before the new one: say the Gujarati for a pearl, then the Gujarati for to mend (a tear).
\end{quote}

\subsection*{You'll want to know: \gu{થીંગડું}}
\textbf{\gu{થીંગડું}} — \emph{thī̃gḍũ} — \textquotedblleft{}a patch (on cloth)\textquotedblright{}.

\begin{quote}
\gu{મારા} \gu{જૂના} \gu{ખમીસ} \gu{પર} \gu{એક} \gu{થીંગડું} \gu{છે}.
\end{quote}

\begin{quote}
\emph{mārā jūnā khamīs par ek thī̃gḍũ che.}
\end{quote}

\begin{quote}
There is a patch on my old shirt.
\end{quote}

A neuter noun.

\begin{cousinweb}
\textbf{\gu{થીંગડું}} (\emph{thī̃gḍũ}) comes from \emph{*thigga}, a form CDIAL reconstructs, and CDIAL lists the Gujarati word as \textquotedblleft{}a patch, a rag\textquotedblright{} (CDIAL 6096).
\end{cousinweb}

\begin{grammarlens}[title={What you've built}]
One more word for describing what you lived and felt.
\end{grammarlens}

\subsection*{Your turn}
\begin{itemize}
\raggedright
  \item \emph{Say it:} \emph{thī̃gḍũ}
  \item \emph{Say it:} \emph{thī̃gḍũ}, once more
  \item \emph{Say it:} \emph{mārā jūnā khamīs par ek thī̃gḍũ che.} — There is a patch on my old shirt.
  \item \emph{Recall:} say the Gujarati for a pearl, then the Gujarati for to mend (a tear), then say \emph{thī̃gḍũ} again
\end{itemize}

\begin{tcolorbox}[breakable,colback=teal!4,colframe=teal!35!black,title={Before you move on}]
What does \gu{થીંગડું} mean? (\textquotedblleft{}A patch (on cloth)\textquotedblright{}.) Say it once more.
\end{tcolorbox}
\section[\texorpdfstring{\gu{ભરતકામ} (bharatkām)}{ભરતકામ (bharatkām)}]{\gu{ભરતકામ} (bharatkām) — embroidery}

\label{lesson:GU-C507-bharatkam}



\begin{quote}
Before the new one: say the Gujarati for to mend (a tear), then the Gujarati for a patch (on cloth).
\end{quote}

\subsection*{You'll want to know: \gu{ભરતકામ}}
\textbf{\gu{ભરતકામ}} — \emph{bharatkām} — \textquotedblleft{}embroidery\textquotedblright{}.

\begin{quote}
\gu{મારી} \gu{માસી} \gu{સુંદર} \gu{ભરતકામ} \gu{કરે} \gu{છે}.
\end{quote}

\begin{quote}
\emph{mārī māsī sundar bharatkām kare che.}
\end{quote}

\begin{quote}
My aunt does beautiful embroidery.
\end{quote}

A neuter noun.

\begin{cousinweb}
\textbf{\gu{ભરતકામ}} (\emph{bharatkām}) is built from \emph{bharat}, \textquotedblleft{}filling in\textquotedblright{}, from \textbf{\gu{ભરવું}}, \textquotedblleft{}to fill\textquotedblright{} (CDIAL 9397), and \textbf{\gu{કામ}}, \textquotedblleft{}work\textquotedblright{}: thread fills in the pattern.
\end{cousinweb}

\begin{grammarlens}[title={What you've built}]
One more word for describing what you lived and felt.
\end{grammarlens}

\subsection*{Your turn}
\begin{itemize}
\raggedright
  \item \emph{Say it:} \emph{bharatkām}
  \item \emph{Say it:} \emph{bharatkām}, once more
  \item \emph{Say it:} \emph{mārī māsī sundar bharatkām kare che.} — My aunt does beautiful embroidery.
  \item \emph{Recall:} say the Gujarati for to mend (a tear), then the Gujarati for a patch (on cloth), then say \emph{bharatkām} again
\end{itemize}

\begin{tcolorbox}[breakable,colback=teal!4,colframe=teal!35!black,title={Before you move on}]
What does \gu{ભરતકામ} mean? (\textquotedblleft{}Embroidery\textquotedblright{}.) Say it once more.
\end{tcolorbox}
